\section{Related Work}\label{sec:related}

\paragraph{ASISR and scale-dependent readout.}
LIIF \citep{chen2021liif} conditions local decoding on a two-dimensional
cell. Its original ablation already varies cell size at fixed $\times30$
queries and discusses area averaging. LTE \citep{lee2022lte} estimates
local Fourier amplitude and frequency, with cell size entering a learned
phase estimator. The examined SC-INR implementation modulates features
with sinc factors and includes separate scale-consistency components.
These constructions are distinct; common band ordering does not prove a
shared integration mechanism. MetaSR \citep{hu2019metasr}, SRNO
\citep{wei2023srno}, CiaoSR \citep{cao2023ciaosr} and HIIF
\citep{jiang2025hiif} further develop continuous decoders, while a
controlled re-training benchmark \citep{nasir2026empirical} examines
architecture and recipe effects.

\paragraph{Integrated encoding and analytic filtering.}
IPE-LIIF \citep{liu2021ipe} integrates sinusoidal positional encodings over
a pixel area, yielding sinc attenuation, and compares this encoding with
plain positional encoding and an additional cell input. Thus neither
area-dependent encoding nor its comparison with learned cell conditioning
is new here. Thera \citep{becker2025thera} uses neural heat fields with
analytic Gaussian filtering for arbitrary-scale readout; Gaussian-splatting
decoders \citep{chen2025gsasr} offer another representation. Integrating
features before a nonlinear RGB decoder is different from a certified
integral of its output, and box and Gaussian target measurements must be
distinguished. Our current evidence measures paired frequency recovery;
it does not establish an advantage over these integral-readout methods.

\paragraph{Frequency diagnostics.}
HyRA \citep{fsds2024} studies impulse responses and a linear/nonlinear
response decomposition, and introduces Frequency Spectrum Distribution
Similarity (FSDS); it is not limited to one global scalar. Our measurement
instead compares signal-normalised reconstruction error in stated bands
and LR Nyquist units. Fourier ring correlation
\citep{vanheel2005frc} and scale-decomposed skill scores are related
agreement diagnostics. We claim neither the first frequency analysis of
SR nor a new information measure.

\paragraph{Spectral control and inverse-problem references.}
SIREN \citep{sitzmann2020siren}, Fourier features
\citep{tancik2020fourier,jacot2018ntk} and FINER \citep{liu2024finer}
control a representation's frequency preference. FreSh
\citep{kania2025fresh}, WINNER \citep{winner2025}, SEC \citep{sec2026}
and WIEN-INR \citep{wieninr} use spectral summaries or multiresolution
components. These signal-fitting and inverse-task studies differ from
our primary pretrained ASISR setting. SuperF \citep{superf2026}, DIP
\citep{ulyanov2018dip} and GP regression \citep{rasmussen2006gp} provide
context for the supplementary inverse checks. Classical LMMSE and
sampling analysis \citep{kay1993estimation} supply conditional linear
references, not a natural-image recovery ceiling. Spectral bias
\citep{rahaman2019spectral} also does not identify a trained decoder's
scale-readout mechanism.
